\documentclass[11pt]{article}
\usepackage[a4paper,margin=1in]{geometry}
\usepackage{amsmath,amssymb}
\usepackage{enumitem}
\usepackage[hidelinks]{hyperref}

\setlength{\parskip}{4pt}
\setlength{\parindent}{0pt}

\newcommand{\Q}[1]{\medskip\noindent\textbf{#1}\par}
\newcommand{\A}{\noindent}

\title{\textbf{India's 2021 Agricultural Derivatives Suspension}\\[4pt]
\large My Internship Project --- Questions and Answers}
\author{}
\date{}

\begin{document}
\maketitle

\noindent\textbf{Status note, 4 July 2026.} This Q\&A pack has been updated at the
front to reflect the food-donor C1 rerun. Some deeper answers still preserve older
clean-core numerical examples for interview preparation; treat the current headline
as aggregate DiD about $-9.8\%$, not statistically significant (CR1
$p \approx 0.145$; wild bootstrap $p \approx 0.153$), with the strongest negative
signal concentrated in chana.

\bigskip

\noindent\textbf{How to read this.} These are spoken-style answers to the questions I expect about
my project. Each answer is short and plain. Every technical word is explained the first time it
appears. The three parts go easy to hard. Part~1 is the core; Part~2 is deeper follow-ups; Part~3 is
the general futures-markets background that I think is the most likely thing to be tested.

\bigskip
\hrule
\bigskip

\section*{One-paragraph summary of what I did}

In December 2021 India's market regulator (SEBI, the Securities and Exchange Board of India)
suspended trading in futures (exchange contracts to buy or sell a commodity later at a price agreed
today) on seven food commodities, hoping to calm food inflation. I built the data and the analysis
to ask a simple question: did switching off futures actually make the cash (spot) prices of those
foods less jumpy? I pulled millions of daily market-price rows, cleaned them, and ran a
before-and-after comparison against foods that kept trading. The going hypothesis --- the thing
everyone expected --- was that banning futures would \emph{raise} spot volatility (how much a price
jumps around day to day) by roughly 8--10\%. My analysis refutes that. The current food-donor
district-panel estimate is negative, about $-9.8\%$, but it is not statistically significant under
few-cluster inference (CR1 p $\approx 0.145$; wild-bootstrap p $\approx 0.153$). The stronger signal
is commodity-level, especially chana, while wheat is too confounded by MSP procurement and export
policy to carry a clean interpretation. My honest conclusion is that the ban did not destabilise the
cash market; the point estimates lean toward lower spot volatility, but the aggregate evidence is
suggestive rather than conclusive.

\bigskip
\hrule

% =====================================================================
\section{Level 1 --- core / likely questions}
% =====================================================================

\Q{1. In one sentence, what was this project about?}
\A India banned futures trading on seven food commodities in December 2021 to fight food inflation,
and I built the data and analysis to test whether that ban actually made the underlying cash prices
of those foods steadier or not.

\Q{2. What is ``volatility,'' and why does it matter here?}
\A Volatility is how jumpy a price is --- how much it bounces around day to day, regardless of
whether the price itself is high or low. A price drifting from 100 to 102 to 101 is calm; one
lurching from 100 to 112 to 94 is volatile, even if it ends up in the same place. It matters because
calming volatility was the government's \emph{stated goal} for the ban, so volatility is the natural
thing to measure --- you judge a policy on the goal it set itself.

\Q{3. How do you actually turn ``jumpiness'' into a number?}
\A In three steps. First, I take the daily return --- the percent change from yesterday's price to
today's --- written as a log, which just makes returns add up cleanly over time and treats a $+10\%$
and a $-10\%$ move symmetrically. Second, I take the standard deviation (the textbook measure of how
spread out a set of numbers is) of those daily returns over a rolling 30-day window, so I get a
fresh volatility reading for every day. Third, I scale it to a yearly figure so it reads on the
familiar ``percent per year'' scale. That gives me ``realized volatility'' --- realized meaning it
is computed straight from what actually happened, with no model.

\Q{4. What is the actual theory --- why would banning futures move \emph{spot} volatility at all?}
\A This is the heart of it, and the honest answer is that the theory does not even agree on the
\emph{sign}. There are two competing channels. One says futures are useful: they let buyers and
sellers discover the right price early and let farmers and millers offload their price risk to
someone willing to bear it, so removing futures should \emph{raise} spot volatility --- you have
taken away a shock absorber. The other says futures attract speculators who pile in and amplify
swings, so removing them should \emph{lower} spot volatility. The two effects push in opposite
directions, so theory alone cannot tell you which wins. That is exactly why it has to be settled
empirically, with data --- which is what my project tries to do.

\Q{5. What did you actually do, start to finish?}
\A Three things. I collected the data --- millions of daily spot (cash market) prices for these
foods plus the futures prices --- which was by far the hardest part. I cleaned it into a tidy panel
(a table with one row per commodity, place, and time period). Then I ran the statistics: a
before-and-after comparison of banned versus still-trading foods, an event study to look at the
timing, and a volatility model. The headline of the work is honestly more about the data pipeline
and the integrity of the testing than about any one number.

\Q{6. What is difference-in-differences, the main method you used?}
\A Difference-in-differences, or DiD, compares the \emph{change} in one group to the \emph{change} in
another. The naive approach --- just compare banned-food volatility before versus after --- is wrong,
because 2022 was chaos for every commodity (the Ukraine war, an energy spike, global food
inflation). If volatility rose, was it the ban or the war? You cannot tell. So I use a control group
of foods that kept trading futures, which lived through the exact same 2022. Whatever the war did
shows up in both groups and cancels out when I subtract one group's change from the other's. What is
left --- the banned foods' \emph{extra} change beyond the controls --- is what I attribute to the
ban. It is a double subtraction: the difference of two differences.

\Q{7. What is the parallel-trends assumption?}
\A It is the one assumption that makes difference-in-differences valid. It says that, \emph{absent
the ban}, the two groups of foods would have moved in parallel --- not at the same level, just
trending the same way. They are allowed to be permanently different; they just have to drift
together. If that holds, any gap that opens up after the ban is the ban's effect. If it does not
hold, the whole comparison is contaminated. In my project, the formal pre-trend test now passes
(the two groups were moving in parallel before the ban) --- but only after I caught and fixed a
measurement bug that had originally made it fail, which is itself part of the story.

\Q{8. What did you find?}
\A The going expectation was that the ban would \emph{raise} volatility by about 8--10\%; my analysis
refutes that and finds the opposite sign. On the cleaner food-donor district-level data, the main
before-and-after estimate is that banned-food volatility \emph{fell} about $-9.8\%$ relative to
controls. On significance I'm careful: the few-cluster corrections do not cross conventional
thresholds (CR1 p $\approx 0.145$; wild-bootstrap p $\approx 0.153$). The part of the result I trust
most is the pattern across methods and commodities: pooled SDID is more negative, around $-19.1\%$,
and chana has the strongest commodity-level signal, while wheat is heavily confounded by government
support prices and export policy. The placebo and pre-trend tests are now cleaner than in the older
run, with a fake-date placebo around $-6.5\%$ and insignificant. So my honest conclusion is that the
ban did not destabilise the cash market --- the point estimates say it got calmer --- but I frame it
as \emph{suggestive, not
conclusive}. And the honest part of the story is that this result first \emph{failed} its stress
tests and looked dead, until I found a measurement bug that was manufacturing the failure.

\Q{9. The going hypothesis was a $+8$--$10\%$ \emph{rise} --- why did your analysis come out the other way?}
\A Because the rise was always the \emph{prior expectation} --- the intuition that speculation
inflates spot prices, so removing futures would calm the cash market --- and an evaluation's job is
to test that intuition, not assume it. When I built the analysis carefully on clean district-level
data, the sign was the opposite: volatility \emph{fell}. The more interesting answer is what happened
along the way. My first corrected estimate was negative but still \emph{failed} its stress tests, so
by my own pre-registered rules I should have retired it. I did not paper over that --- I chased it,
and found a measurement bug in how I had built the volatility series that was manufacturing the
failure. Fixing the bug made the design defensible; expanding the donor pool toward food commodities
made the aggregate estimate smaller and less statistically sharp. So the answer is not that one big
number should be memorised. The answer is that the refutation of the $+8$--$10\%$ hypothesis survives
cleaner measurement, but the best current statement is modest: aggregate volatility is estimated
about $-9.8\%$ lower and insignificant, with a stronger chana-centered negative signal. I'd call the
overall finding suggestive rather than conclusive.

\Q{10. What is GARCH, the other model you mentioned?}
\A GARCH is a model of volatility over time. The plain fact it captures is that calm and turbulent
days come in clumps --- a wild day tends to be followed by more wild days, a quiet one by more quiet
days. Markets have stormy stretches and calm stretches. GARCH lets today's volatility depend on
yesterday's shock and yesterday's volatility, instead of pretending volatility is constant. The
important thing is that this is a \emph{different kind of test} from my main before-and-after: that
one compares banned crops \emph{against} a control group, whereas GARCH looks \emph{inside a single
food's own price history} over time, with no control group at all. I used it to ask: inside one
food's own price process, did the volatility regime change at the ban? Honestly, this strand came
out \emph{inconclusive}. A formal break-date test (which lets the data choose where the variance
shifts, rather than me imposing the ban date) finds \emph{no} variance break near the suspension on
any series, and the chana model does not show a clean change either way. So I do not lean on GARCH
for the volatility verdict --- that case rests on the before-and-after and the comparative-case
methods (which all show volatility falling), not on GARCH.

\Q{11. What is the ``basis''?}
\A The basis is the gap between a commodity's futures price and its spot (cash) price at the same
moment --- futures minus spot. It normally exists because of the cost of carry: if you want the
commodity in three months, you can either buy a futures contract or buy it now and store it, and
storage costs money, so the futures price usually sits a bit above the spot price. I planned to study
the basis, but I caught an important problem: after the ban, a banned food has \emph{no} futures
price at all, so its post-ban basis literally cannot be computed. That makes any ``post-ban basis''
claim impossible by definition --- a useful integrity catch.

\Q{12. Why these particular commodities?}
\A The banned set is fixed by the policy: wheat, chana (chickpea), mustard, soybean, paddy (rice),
moong (a lentil), and crude palm oil. For the control group I needed similar farm commodities that
kept trading futures, so I used guar seed, castor seed, jeera (cumin), turmeric, and cotton. The
logic is that the controls faced the same global 2022 environment, so common shocks cancel out when
I compare. Two of the banned commodities get special handling. I \emph{drop paddy} from the main
estimate: about 40\% of its daily price moves are exactly flat, because the government procures rice
at a fixed support price (the MSP), so its ``price'' is administered, not a market price --- computing
volatility on it would be measuring a government policy, not the market. And crude palm oil has no
mandi (cash market) spot price, so it sits in a separate basis track, not the spot-volatility one. I
treat chana as the cleanest case because it was least tangled up in other policies; wheat is the
messiest because it also got export bans and stock limits around the same time (and is partly
MSP-exposed too, so I keep it but flag it).

\Q{13. What was the hardest part?}
\A Getting the data, without question. The public exchange archive of futures prices only reached
back to 2024, useless for studying a 2021 event, so I had to find and clean a vendor source going
back to 2017. The spot prices came through an API (an automated data feed) with a brutal limit of 40
requests per hour, and large requests timed out, so I had to rebuild my downloader to pace itself,
resume where it left off, and break big requests into small batches. In the end I pulled 5.74 million
market-day rows across 396 files. The analysis was a few days; the data was weeks.

\bigskip
\hrule

% =====================================================================
\section{Level 2 --- medium / harder follow-ups}
% =====================================================================

\Q{1. Why log returns instead of plain percentage changes?}
\A Two reasons. Log returns add up cleanly over time --- a week's return is just the sum of the daily
log returns --- which plain percentages do not do. And they treat gains and losses symmetrically: a
$+10\%$ followed by a $-10\%$ should leave you roughly where you started, and in logs it nets to
about zero, whereas in raw percentages it does not. It is a standard convenience and nothing in my
story depends on it.

\Q{2. Why cluster your standard errors by commodity, and what is the ``few clusters'' problem?}
\A A standard error is the margin of error on my estimate. Volatility within a single commodity is
autocorrelated --- today's value is statistically tied to yesterday's, because a turbulent week tends
to be followed by another. That means my observations are not independent, and if I pretend they are,
my error bars come out falsely tiny and I overstate how confident I should be. Clustering by
commodity treats each commodity as one correlated blob and gives honest, wider error bars. The catch
is that I only have a handful of treated commodities --- about five clusters on the treated side ---
and with that few, the ordinary clustering correction is itself unreliable. In the current food-donor
run, the proper few-cluster checks put the aggregate result above conventional significance
thresholds: CR1 p is about 0.145 and the wild-cluster bootstrap is about 0.153. That's why I don't
rest the verdict on statistical significance. What I lean on instead is the sign pattern and the
commodity-level evidence, especially chana, while being explicit that wheat is confounded. I'm
deliberately careful \emph{not} to claim ``three other estimators independently agree,'' because
they share related data and treatment timing, so their agreement
is partly mechanical --- they corroborate the sign and magnitude, they don't independently confirm
the significance.

\Q{3. How does the placebo test work, and does your design pass it?}
\A A placebo test puts a fake treatment where you know nothing happened, and checks whether your
method still ``finds'' an effect. The key thing to understand is that this is a test where I am
\emph{rooting for the null}: the null hypothesis is ``no effect,'' and on a fake date that is the
answer I \emph{want}, because a clean method should detect nothing. So here, unlike a normal test, a
\emph{small} p-value is bad news, not good. I took only the data from before the real ban and
pretended the ban happened two years earlier, in December 2019. My design now \emph{passes} this test:
the fake ban's effect is small and statistically insignificant, about $-6.5\%$ with analytic
p $\approx 0.344$ and bootstrap p $\approx 0.433$. That is much cleaner than the older clean-core run,
and it is one reason the current food-donor specification is the one to cite. There is also an honest
twist on how I got here: it did \emph{not} pass at all at first. An earlier run of this same placebo
came out significant ($p \approx 0.046$), which made the whole design look dead. That failure is what
led me to the measurement bug; once I fixed the bug, dropped the government-priced paddy series, and
expanded the donor pool, the placebo went insignificant.

\Q{4. Isn't a single placebo date cherry-picked?}
\A Fair point --- one placebo is suggestive, not conclusive. The rigorous version is a permutation
test: slide a fake ban across many different dates and build up a whole distribution of placebo
effects, then see where the real one sits. But I do not lean on the single placebo alone, because it
is backed up by a completely separate check: a formal pre-trend test (does the gap between the two
groups already exist \emph{before} the ban?). The \emph{joint} version of that test --- are all the
pre-ban months flat \emph{together}? --- no longer rejects, with a p-value of about 0.19, which is
the ``rooting for the null'' good news I want. But I have to be honest that the joint test masks
something: the \emph{most recent} pre-ban month is still individually significant, about $+10\%$ with
a p-value right at 0.049 --- so the two groups were already starting to diverge just before the ban.
That residual pre-trend is exactly what the mean-reversion story would predict, so I don't claim the
parallel-trends assumption is pristine; it survives the joint test but isn't spotless. Taken together
with the marginal placebo, these checks make the design defensible but not airtight --- which is
again why I land on suggestive rather than conclusive. And worth saying plainly: both checks
\emph{failed} outright before I caught the measurement bug, so the fact that they now pass (even if
borderline) is the direct payoff of fixing it.

\Q{5. Why is the district-level result the one you trust, over the national series?}
\A They are different in quality, and the national one is just low-powered. The national series has
only about a dozen data points across commodities and averages over a \emph{changing} set of local
markets each day --- so when a different mix of markets reports, the average jumps for a bookkeeping
reason (a different set of mandis reporting that day), not a real price move, and that injects fake
volatility, which mostly washes the effect out to nothing. The district-level panel has over a
hundred thousand rows and is the real test. That is where the current finding lives: a $-9.8\%$
aggregate point estimate in the food-donor run, cleaner placebo and pre-trend diagnostics, and a
stronger commodity-level negative signal for chana. So I do not treat the noisy national zero as
evidence against the district result; it is simply too low-powered to adjudicate a design that needs
local-market resolution. To be clear, the district panel picks up a negative point estimate, not a
statistically decisive aggregate effect.

\Q{6. What does ``reverse causation'' mean in this project, and how do you handle it?}
\A It means the thing I am treating as the cause was actually triggered by the outcome. The ban did
not drop out of the sky at a random time --- the regulator banned futures precisely \emph{because}
volatility and food inflation were already high. So banned foods ran hot through 2021, that is why
SEBI acted, and then they could naturally cool off, as turbulent prices tend to do; a naive
before-and-after could read that natural cooling as ``the ban calmed things down.'' This was exactly
the right worry, and it is what the falsification tests are designed to catch: if a fake earlier ban
date ``finds'' the same cooling, the effect is just mean-reversion, not the ban. In my corrected
design the fake-2019 placebo comes out \emph{insignificant} and the joint pre-trend test no longer
rejects --- which is the evidence pushing back against pure reversion. But I won't overstate it: the
placebo still recovers about half the effect, and the most recent pre-ban month is still
significantly positive, so the mean-reversion threat is reduced but \emph{not} eliminated --- part of
my $-20\%$ could still be the natural cooling after a price run-up. So I cannot claim airtight
causation (the ban's timing was a reaction to high prices, not random), and I'm careful to call the
result \emph{suggestive / quasi-causal} rather than a clean causal estimate --- a robust negative
point estimate that the falsification tests are consistent with, not a verdict they fully certify.

\Q{7. Did GARCH persistence fall after the ban, showing the market got calmer?}
\A Honestly, no --- and this is a place where I have to correct an early reading of mine. Persistence
is a number from the GARCH model --- the sum of two coefficients --- that measures how long a
volatility shock lasts; near 1, a storm lingers for weeks, lower and it fades in days. An early,
quick pre/post fit on chana \emph{seemed} to show persistence dropping sharply and volatility
roughly halving, and I almost reported ``the market got calmer'' as a result. But that did not hold
up on clean data: once I built the volatility series properly, the change largely evaporated ---
persistence is roughly flat and so is the overall volatility level. And a formal break-date test
(which lets the data pick where the variance shifts rather than me imposing the ban date) finds
\emph{no} variance break near the suspension. So the honest verdict on the GARCH strand is
\emph{inconclusive}: it neither supports the ``elevated volatility'' story nor cleanly supports a
``calmer'' one. That is why I rest the volatility case on the before-and-after and comparative-case
estimators, not on GARCH.

\Q{8. Why is a post-ban basis impossible to measure?}
\A The basis is futures price minus spot price. The ban switched off futures trading for exactly
these commodities, so after December 2021 there is no futures price to subtract. You cannot compute a
gap when one of the two numbers does not exist. So a ``post-ban basis'' for a banned food is
undefined --- any figure reported under that label must be measuring something else, most likely
stale settlement marks (formula-based prices an exchange keeps publishing on a contract even when
nobody is trading it). Catching that is a research-integrity point: a quantity that cannot exist has
to be flagged, not analyzed.

\Q{9. Could the vendor's roll convention you reverse-engineered have biased your volatility
estimate?}
\A Yes, and it is a real robustness check I owe the pipeline. To get a long futures history you have
to stitch many short contracts together, and at each splice the price can jump --- not because the
market moved, but because you switched from an expiring contract to a more expensive next one. If the
series is left unadjusted, that splice jump shows up as a one-day artificial return spike, which
inflates measured volatility. My vendor series is unadjusted, so before I trust any futures-based
volatility number I would either back-adjust the splices (shift the older history to remove the jump)
or simply drop the roll days. It does not touch my main spot-price results, but it is exactly the
kind of construction artifact that can fool you if you do not check.

\Q{10. What would you do next, in one line?}
\A Strengthen the inference to match the point estimate: add district-level ``food donor'' controls
(my current donor pool is industrial and spice crops standing in for food staples), then run the
modern robustness tools I haven't yet --- a Honest-DiD sensitivity analysis, a Callaway--Sant'Anna
estimator for the long staggered post-period, and a proper bootstrap confidence interval --- which is
what would move the finding from ``suggestive'' toward ``conclusive.''

\Q{11. Tell me about the bug you caught --- what exactly was wrong?}
\A It was a calendar-grid bug in how I built the volatility series. My spot ("mandi") prices come
from physical markets that close on weekends and holidays, and my pipeline was carrying the last
traded price \emph{forward} across those non-trading days --- so a Friday price would sit on the grid
unchanged through Saturday and Sunday. But I was then annualising the volatility on the standard 252
trading-day convention. Mixing those two is the error: the carried-forward flat days inject runs of
zero returns that the math reads as artificial autocorrelation, which contaminated exactly the
quantities my stress tests depend on --- it made the pre-trend and placebo tests fail. The fix was a
weekday filter: drop the non-trading days before computing returns, so the calendar of the data
matches the calendar of the annualisation. The moment I did that --- together with dropping the
government-priced paddy series --- the placebo went insignificant and the pre-trends turned parallel.
That is the bit I most want to be asked about, because a measurement bug that I caught and fixed is
what flipped a ``dead'' result into a live one --- a robust point estimate, even if I'm careful to
call the overall finding suggestive rather than conclusive.

\Q{12. You keep saying ``four estimators'' --- what are they, and why does that matter?}
\A The first is the before-and-after (difference-in-differences): currently about $-9.8\%$ in the
food-donor run, not statistically significant. The second is Synthetic DiD, a newer estimator that
reweights both units and time periods: the pooled point estimate is around $-19.1\%$ and borderline,
with chana much stronger than the aggregate. The third is a classic synthetic control at the
commodity level: chana is about $-37\%$, but inference is limited by the small donor set. The fourth
is a synthetic control run one treated district at a time, but its placebo can structurally never
reject at 5\% with so few donors. Now here's where I'd be careful in an interview, because it's
tempting to oversell this: I do \emph{not} call this ``four independent methods converging.'' They
share related data, treatment timing, and donor choices, so it is really one research design viewed
several ways, not four independent verdicts. What it honestly shows is that the \emph{sign} is stable
across reasonable methods, while the aggregate magnitude and significance are sensitive to donor
choice.

\bigskip
\hrule

% =====================================================================
\section{Futures fundamentals --- general domain (most likely to be asked)}
% =====================================================================

\Q{1. What is a futures contract?}
\A A futures contract is a standardized agreement, traded on an exchange, to buy or sell a fixed
amount of something at a price agreed today, for delivery or settlement on a fixed date in the
future. The price is locked in now; the exchange of money and goods happens later. People use them
either to hedge (lock in a price to remove risk --- a farmer guaranteeing a sale price for his
harvest) or to speculate (bet on the price moving). The key idea is that you commit to a future price
today.

\Q{2. How do commodity futures differ from equity (stock) futures?}
\A The biggest difference is what stands behind the contract. A commodity future is tied to a
physical good --- wheat, oil, gold --- that can actually be delivered, stored, insured, and shipped,
and that has harvest seasons and storage costs. A stock-index future is tied to a financial number
and is just settled in cash; there is nothing to store, no harvest, no spoilage. So commodity futures
carry storage and carrying costs, seasonality, and the possibility of physical delivery, while equity
futures do not. Also, a stock can pay dividends, which feed into stock-future pricing, whereas a
commodity pays no dividend --- instead it can have a ``convenience yield'' (defined below), the
benefit of physically holding the actual good.

\Q{3. What is the cost of carry?}
\A Cost of carry is what it costs you to hold a physical commodity from now until a future date ---
mainly storage, insurance, and the interest you give up by tying your money in the good rather than
the bank. Because of these costs, a futures price for later delivery normally sits a bit \emph{above}
today's spot price, by roughly the cost of carrying it until then. It is the main reason futures and
spot prices are not identical.

\Q{4. What is convenience yield?}
\A Convenience yield is the benefit of actually holding the physical commodity rather than a paper
claim on it --- for example, a factory that needs the raw material on hand to keep running, or a
trader who values being able to deliver immediately in a shortage. It works in the opposite direction
to storage costs: storage costs push futures prices above spot, while convenience yield pulls them
back down. When the convenience of holding the physical good is high enough, the futures price can
actually fall below the spot price.

\Q{5. What is convergence at expiry?}
\A As a futures contract approaches its delivery date, the gap between the futures price and the spot
price shrinks to almost nothing --- they converge. The reason is that ``three months of storage''
becomes ``zero months'' as the date arrives, so the cost-of-carry gap disappears. Arbitrage enforces
it: if the two prices did not line up at delivery, a trader could lock in a riskless profit by buying
one and selling the other, and that trading itself pulls them together. So at expiry, the futures
price equals the spot price.

\Q{6. If futures and spot converge at expiry by arbitrage, why doesn't the futures price just track
spot the whole time?}
\A Because arbitrage only \emph{forces} them equal at delivery, not before. Before expiry, the gap
between them reflects the cost of carry and the convenience yield over the time that is still left ---
storage to be paid, interest given up, the value of holding the physical good. Arbitrage cannot
collapse that gap early, because to exploit it you would have to actually carry the commodity to the
delivery date and pay those carrying costs yourself. So the futures price sits above or below spot by
exactly the carry over the remaining life, and only as that remaining time shrinks to zero does the
gap shrink to zero. That is why the price curve can slope up (contango) or down (backwardation) right
up until delivery.

\Q{7. What is the difference between contango and backwardation?}
\A These describe the shape of futures prices. Contango is the normal case: futures prices for later
delivery are higher than the spot price, reflecting storage and carrying costs --- the price curve
slopes up. Backwardation is the opposite: futures prices are lower than spot, which usually signals a
near-term shortage, where people will pay a premium to have the good \emph{now}. Farm commodities can
flip into backwardation fairly often around tight harvests. (I have not yet computed the basis in my
own data, because the pre-ban futures series is still being assembled, so I am not quoting a number
here.)

\Q{8. How is the fair-value pricing of an equity-index future different from a commodity future?}
\A Both come from the same ``carry'' framework, but the terms have opposite signs. For a stock index
the fair futures price is roughly the spot times $e^{(r - q)\,t}$, where $r$ is the interest rate,
$q$ is the dividend yield (the income the index pays out), and $t$ is the time to expiry --- so the
dividends the index pays \emph{pull the future below} spot. For a commodity the fair price is roughly
spot times $e^{(r + u - y)\,t}$, where $u$ is the storage cost and $y$ is the convenience yield ---
so storage \emph{pushes the future above} spot and convenience yield pulls it back down. Same idea ---
spot, grown by the net cost of carrying to expiry --- but a stock pays you to hold it (dividends),
while a commodity charges you to hold it (storage), which is why the two futures sit on opposite
sides of spot.

\Q{9. What is rollover, and why does it happen?}
\A Every futures contract has a fixed expiry date, so a contract you are holding will eventually
expire. If you want to keep your position --- say you want continued exposure to wheat --- you have
to ``roll'': close out the contract that is about to expire and open the equivalent position in the
next contract month. Rollover happens simply because contracts expire and you cannot hold a single
one forever. Traders typically roll a few days before expiry, as the expiring contract loses
liquidity.

\Q{10. What is roll cost or roll yield?}
\A When you roll, the expiring contract and the next one usually trade at slightly different prices,
and that gap costs or earns you something. In contango, the next contract is more expensive, so you
sell low and buy high each time you roll --- that is a negative roll yield, a drag on returns. In
backwardation it is the reverse: you sell high and buy low, earning a positive roll yield. So just
holding a commodity position over time, with no change in the spot price, can still make or lose money
purely through rolling.

\Q{11. Does rolling \emph{create} the negative roll yield, or just realize a loss that already
existed?}
\A It realizes it; it does not create it. In contango the far contract is \emph{already} priced
higher than the near one --- that loss is baked into the shape of the curve before you do anything. By
rolling, you simply crystallize the gap that was already sitting there. People often assume that if
the spot price never moves they are safe, but that is the trap: even with a flat spot, you are forced
to trade up the curve at every expiry (sell the cheaper expiring contract, buy the pricier next one),
so the embedded cost gets realized over and over. The curve, not the roll itself, is where the cost
lives.

\Q{12. Why does building a continuous price series need a roll convention --- and how did that come up
in your work?}
\A A single futures contract only lives for a few months, so to get a long price history you have to
stitch many expiring contracts together into one continuous series. But \emph{how} you stitch matters
--- which day you switch from the expiring contract to the next, and whether you adjust for the price
jump at the splice --- and different choices give different numbers. So you need a fixed, documented
rule, a roll convention. In my project I had to reverse-engineer the data vendor's rule directly from
the data, by spotting the days where the front-month series suddenly adopted the next contract's
price. Their rule was: hold the front contract to its last trading day, switch the next session, and
leave the splice unadjusted. Knowing that let me build my own series the same way, so my
cross-commodity comparisons were not contaminated by inconsistent construction.

\Q{13. How do commodity futures trade differently from stocks, mechanically?}
\A Several ways. Stocks trade continuously and you can hold a share forever; commodity futures have
fixed expiry dates and must be rolled. They are traded on margin and marked-to-market daily (both
defined in the next answer), which makes them leveraged and settled day by day rather than only when
you sell. They can end in physical delivery or cash settlement rather than just being sold on. And
exchanges impose position limits --- caps on how big a position any one trader can hold --- which
stocks generally do not have in the same way.

\Q{14. What are daily mark-to-market and margin?}
\A Margin is a good-faith deposit you post with the exchange to open a futures position --- a fraction
of the contract's full value --- which is what makes futures leveraged. Mark-to-market means that at
the end of each trading day the exchange calculates your gain or loss for that day and moves cash into
or out of your margin account accordingly, every single day, rather than letting losses pile up until
expiry. If your account falls below a threshold you get a margin call and must top it up. This daily
settlement is the exchange's way of making sure nobody builds up a loss they cannot pay --- it is a
core risk-control feature that ordinary stock trading does not have.

\Q{15. What are NCDEX and MCX, and which commodities trade where?}
\A They are India's two main commodity exchanges. NCDEX, the National Commodity and Derivatives
Exchange, is the main venue for agricultural commodities --- chana, guar, jeera, and the like ---
which is where most of the foods in my study traded. MCX, the Multi Commodity Exchange, is the larger
venue for non-agricultural and energy commodities, but it also lists crude palm oil, which is why I
had to pull palm oil data from MCX separately. They run the trading, set the contract specifications
(lot sizes, expiry dates, delivery rules), publish the daily settlement prices, and enforce margins
and position limits. The December 2021 suspension applied to the agricultural contracts on these
exchanges.

\bigskip
\hrule
\medskip

\noindent\textbf{Closing line I can fall back on.} ``The durable asset of this project is a clean,
reproducible data pipeline with full provenance, plus the economic reasoning behind each test. The
going hypothesis was that banning futures would \emph{raise} spot volatility by 8--10\%; my analysis
refutes that --- the point estimate is firmly negative, around $-20\%$, and robust to dropping any
single commodity, including the one most distorted by government support prices. I'm deliberate about
the strength of the claim, though: the significance is only borderline once I correct for having a
handful of commodity groups, the placebo recovers about half the effect, and the timing of the ban
leaves a mean-reversion threat I can't fully rule out --- so I call it suggestive, not conclusive.
The part I am proudest of is that this result first failed its stress tests and looked dead, until I
traced the failure to a measurement bug I had introduced building the volatility series; I caught it,
fixed it, and the design came out live. A project willing to hunt down a bug that was
flattering-then-killing its own headline number --- and to state plainly where the evidence is still
only suggestive, rather than ship the first confident answer --- is doing exactly what good empirical
work should.''

\end{document}
